\subsection{حمل الرسائل الكبيرة باستخدام تعليقات متعددة}
\label{sec:multiframe}
% Generated by scripts/build_multicomment_example.py.
\newcommand{\McPayloadBytes}{40}
\newcommand{\McPayloadStorageBits}{320}
\newcommand{\McCompressedBits}{31}
\newcommand{\McCount}{7}
\newcommand{\McMarkerBits}{16}
\newcommand{\McVersionBits}{2}
\newcommand{\McCountBits}{5}
\newcommand{\McLengthBits}{9}
\newcommand{\McControlBits}{32}
\newcommand{\McIntegrityBits}{32}
\newcommand{\McStreamBits}{95}
\newcommand{\McPaddingBits}{10}
\newcommand{\McSlotBits}{105}
\newcommand{\McCommentChoices}{1329}
\newcommand{\McTopicChoices}{32}
\newcommand{\McCommentBits}{10}
\newcommand{\McTopicBits}{5}
\newcommand{\McFrameBits}{15}
\newcommand{\McDictionaryEntries}{998}
\newcommand{\McReplyIndex}{662}
\newcommand{\McTopicIndex}{18}
\newcommand{\McFirstReplyBits}{1010010110}
\newcommand{\McFirstTopicBits}{10010}
\newcommand{\McMarkerString}{1010010110100101}
\newcommand{\McVersionString}{10}
\newcommand{\McGammaCount}{00111}
\newcommand{\McGammaLength}{000011111}
\newcommand{\McFrameRowsEn}{%
1 & 662 & 18 & 15 & 0 \\ \hline
2 & 796 & 7 & 15 & 0 \\ \hline
3 & 960 & 2 & 15 & 0 \\ \hline
4 & 0 & 5 & 15 & 0 \\ \hline
5 & 105 & 17 & 15 & 0 \\ \hline
6 & 663 & 5 & 15 & 0 \\ \hline
7 & 288 & 0 & 5 & 10 \\ \hline
}
\newcommand{\McFrameRowsAr}{%
\NUM{1} & \NUM{662} & \NUM{18} & \NUM{15} & \NUM{0} \\ \hline
\NUM{2} & \NUM{796} & \NUM{7} & \NUM{15} & \NUM{0} \\ \hline
\NUM{3} & \NUM{960} & \NUM{2} & \NUM{15} & \NUM{0} \\ \hline
\NUM{4} & \NUM{0} & \NUM{5} & \NUM{15} & \NUM{0} \\ \hline
\NUM{5} & \NUM{105} & \NUM{17} & \NUM{15} & \NUM{0} \\ \hline
\NUM{6} & \NUM{663} & \NUM{5} & \NUM{15} & \NUM{0} \\ \hline
\NUM{7} & \NUM{288} & \NUM{0} & \NUM{5} & \NUM{10} \\ \hline
}


\subsubsection{رسالة واحدة تحملها عدة تعليقات}
يحمل التعليق الواحد اختيارات موضع الرد وموضوع التعليق التي تتيحها القوائم المشتركة. ولا تضيف إطالة نصه اختيارات جديدة إلى هذين الحقلين. وعندما تتجاوز الرسالة ميزانية الاختيار، يستطيع المُرسِل توزيع تتابع بتاتها على عدة تعليقات. نسمي ذلك نمط التعليقات المتعددة \EN{multi-comment}، ونسمي كل تعليق حامل إطارًا من الرسالة نفسها. ولا يلزم أن يكون كل تعليق ردًا على التعليق الحامل السابق؛ إذ تحدد بتات كل إطار موضع الرد الخاص به.

نميّز بين نمط التعليق الواحد ونمط نقل الرسالة عبر تعليقات متعددة. لا تُضاف ترويسة جمع الرسالة الموصوفة أدناه تلقائيًا إلى مسار الاختيار بتعليق واحد، وقد تكون لإعدادات الضغط أو الحماية حقولها الخاصة. أما نمط التعليقات المتعددة، فيضيف المعلومات اللازمة لقراءة التتابع بوصفه رسالة كاملة. وتتبع هذه الكلفة نمط النقل المختار، حتى إذا اتسع تعليق واحد لتتابع النقل كله.

هذا المسار موجود في البرنامج. واستخدمت المقارنات مع \EN{ZLG}~\cite{lin2024zero} مسار التعليق الواحد مع إعادة بناء تستعين بسجل المُرسِل؛ فلا تقيس تلك المقارنات نقل الرسائل الأكبر عبر هذا النمط. وتختبر تجارب مقتطفات المصادر وتجارب الأرشيف شروط استقبال مختلفة للتعليقات المتعددة، يبينها فصل التقييم.

\subsubsection{تحضير الرسالة والتخطيط للتعليقات}
يتفق الطرفان على السياق وقوائم الموضوعات وتمثيل الحمولة ونمط حمايتها. ويحضّر المُرسِل الحمولة كاملة قبل تقسيمها بين التعليقات. يستخدم مسار الحمولة غير المحمية قاموسًا حتميًا من النصوص المشتركة للضغط دون فقد. وتشمل التمثيلات المتاحة \EN{UTF-8} ومؤشرات إلى مقاطع نصية وتكييف شجرة البادئات المستمدة من \EN{OPPRESSION}~\cite{zillien2024oppression}. ويحدد تمثيل الضغط كيفية قراءة الحمولة، بصورة منفصلة عن بادئة النقل المستخدمة لجمع التعليقات.

لتكن $P$ تتابع بتات الحمولة المحضّرة، و$L=|P|$ طوله، و$N$ عدد التعليقات الحاملة. يبني التنفيذ تتابع النقل بالشكل الآتي:
\[
S=M\Vert V\Vert\Gamma(N)\Vert\Gamma(L)\Vert P\Vert R.
\]
تمثل $M$ بادئة تعريف ثابتة، و$V$ إصدار النقل، و$\Gamma$ ترميز عدد صحيح موجب بطريقة \EN{Elias-gamma}، و$R$ فحص \EN{CRC-32} للحمولة $P$. يتيح ترميز العدد تحديد نهايته من البتات نفسها، فيقرأ المُستقبِل العدد والطول دون معرفتهما مسبقًا. ويبلغ طول الحقل:
\[
\ell_\gamma(x)=2\lfloor\log_2 x\rfloor+1.
\]

يحسب المُخطِّط كلفة التتابع كله، بما فيه هذه الحقول. فإذا كانت ميزانية الاختيار القابلة للاستخراج للتعليق $i$ هي $b_i$، وفق القسم~\ref{sec:channel-widths}، فإنه يختار أول عدد من التعليقات يحقق الشرط:
\[
\sum_{i=1}^{N}b_i\ \ge\ L+16+2+\ell_\gamma(N)+\ell_\gamma(L)+32.
\]
ويعيد حساب حقل العدد كلما غيّر القيمة المرشحة لـ$N$. وإذا لم تكفِ المواضع المتاحة لحمل التتابع كاملًا ضمن الحد المسموح، يفشل التخطيط قبل التعامل مع جزء من الرسالة بوصفه رسالة قابلة للإرسال. وفي خيار بناء الموضوعات من سلسلة الرد، يحدد المُخطِّط التعليق الأب أولًا، ثم يعيد حساب ميزانية الموضوعات من سياقه.

يحد عدد التعليقات المسموح من عمل الإرسال وعدد التعليقات الظاهرة، ولا يثبت قيمة $N$ أو يغيّر عدد البتات التي يحملها كل اختيار. وقد يتيح رفع الحد حمل تتابع أطول، مقابل توليد تعليقات أكثر وقراءتها. والحدان المستخدمان في التقييم، وهما ثمانية تعليقات وعشرون تعليقًا، ميزانيتان لتجارب الإرسال، لا شرطان للطريقة أو حجمان أمثلان مثبتان للرسالة.

\subsubsection{وظيفة حقول النقل وكلفتها}
\label{sec:multicomment-overhead}
يفصل الجدول~\ref{tab:multicomment-fields} معلومات جمع الرسالة عن محتواها. وتُكتب هذه الحقول مرة واحدة للرسالة كاملة، ولا تُكرَّر في كل تعليق. وقد تنقسم الحقول نفسها بين تعليقين، مثل بتات الحمولة.

\begin{table}[!htbp]
\centering\small
\caption{حقول تتابع النقل في نمط التعليقات المتعددة، قبل إكمال التعليق الأخير بالأصفار. يعبّر $L$ عن عدد بتات الحمولة المحضّرة، و$N$ عن عدد التعليقات الحاملة.}
\label{tab:multicomment-fields}
\begin{RTL}
\begin{tabularx}{\columnwidth}{|A{0.23\columnwidth}|C|B|}
\hline
الحقل & عدد البتات & الوظيفة \\ \hline
بادئة التعريف $M$ & \NUM{16} & التعرف على صيغة النقل ذات الإصدار. \\ \hline
الإصدار $V$ & \NUM{2} & تحديد إصدار النقل المستخدم. \\ \hline
عدد التعليقات & $\ell_\gamma(N)$ & مقارنة العدد المتوقع بالتعليقات المزودة للمُستقبِل. \\ \hline
طول الحمولة & $\ell_\gamma(L)$ & تحديد نهاية $P$ وفصلها عن الفحص وأصفار الإكمال. \\ \hline
الحمولة المحضّرة $P$ & $L$ & حمل تمثيل الرسالة بعد الضغط أو الحماية. \\ \hline
فحص السلامة $R$ & \NUM{32} & فحص الحمولة المعاد جمعها قبل فك الضغط. \\ \hline
\end{tabularx}
\end{RTL}
\end{table}

بادئة التعريف الثابتة هي تتابع بتات يتفق المُرسِل والمُستقبِل على وضعه في بداية تتابع النقل الكامل. وهي هنا \EN{\texttt{\McMarkerString}}. يضع المُرسِل هذه البتات الستة عشر نفسها في بداية كل رسالة تستخدم نمط النقل هذا، مهما تغير نص الرسالة أو طولها أو عدد التعليقات الحاملة لها. وبعد استخراج بتات التعليقات وجمعها، يقارن المُستقبِل البتات الستة عشر الأولى بالتتابع المتفق عليه. تؤدي البادئة وظيفة علامة معناها: «يستخدم هذا التتابع صيغة النقل متعددة التعليقات ذات الإصدار». وتحملها اختيارات موضع الرد والموضوع مثل سائر بتات التتابع، ولا تُكتب حرفيًا في نصوص التعليقات.

يحدد البتان التاليان الإصدار، وقيمتهما الحالية \EN{\texttt{\McVersionString}}. وعندما تتطابق البادئة والإصدار مع القيم المتوقعة، يقرأ المُستقبِل الحقول التالية بوصفها عدد التعليقات وطول الحمولة والحمولة وفحص \EN{CRC}. وهكذا تعرّف البادئة بصيغة النقل، بينما يحدد حقل الإصدار نسختها. ولا يحدد أي منهما تمثيل الضغط أو يتحقق من صحة الرسالة؛ فهاتان وظيفتا تمثيل الحمولة وفحص السلامة. وتقرأ الشيفرة أيضًا صيغة سابقة تبدأ بالعدد والطول دون هذه البادئة أو فحص \EN{CRC}، وتُوسَم نتيجتها بأنها لم تخضع للتحقق من السلامة. أما \EN{UTF-8} ومؤشرات المقاطع النصية ومؤشرات شجرة البادئات، فهي تمثيلات للحمولة، منفصلة عن التعرف على صيغة النقل.

اختيار التتابع الثابت وطوله قرار في التنفيذ للتعرف على صيغة النقل مع إبقاء مسار قراءة الصيغة السابقة. ولا تثبت السجلات أن \NUM{16} بتًا هي أقل طول ممكن أو أكثره كفاءة. ويمكن لبروتوكول متفق عليه مسبقًا استخدام بادئة أقصر أو الاستغناء عن التعرف على الصيغة، مع تعديل متوافق لدى الطرفين. كذلك يمثل حقل العدد فحص اتساق إضافيًا عندما تكون قائمة التعليقات المرتبة مزودة أصلًا؛ فهو يكشف اختلاف التتابع عن العدد المدوّن، ولا يكتشف التعليقات أو يرتبها تلقائيًا. ويمكن لاتفاق مسبق على طول الحمولة أن يغني عن إرساله في نمط آخر. لذلك يحتاج جمع الرسالة إلى اتفاق صريح على الحدود والترتيب، بينما تبقى عروض الحقول هنا خصائص لهذا التنفيذ.

يتيح \EN{CRC-32} فحص الحمولة المستخرجة، بما في ذلك الحالات التي تنتج فيها بتات خاطئة عن تحديد موضوع غير مقصود. ويحسبه التنفيذ على التمثيل النصي للحمولة $P$ باستخدام المحرفين صفر وواحد. ويعيد المُستقبِل الحساب قبل فك الضغط، ويرفض إعادة البناء إذا اختلفت النتيجتان. ولا يضمن تطابق قيمتي الفحص تطابق الحمولة، فقد تنتج القيمة نفسها عن حمولتين مختلفتين. كذلك لا يصحح الفحص الأخطاء ولا يوفر التوثيق بمفتاح أو سرية المحتوى. ويظل منفصلًا عن نمط حماية الحمولة.

\subsubsection{توليد التعليقات وإكمال التعليق الأخير}
بعد التخطيط، يأخذ المُرسِل مقاطع متتالية من $S$. ويحدد كل مقطع موضع الرد وموضوعه، ثم يولّد نموذج الكتابة تعليقًا لهذين الاختيارين. يُطبّق تحقق المُرسِل على كل تعليق حامل، وتُعاد المحاولة عند رفض المرشح ضمن الحد المتفق عليه. وإذا فشل التوليد، فلا تُعد التعليقات التي أُنتجت قبل الفشل رسالة ناجحة أقصر؛ إذ يتعلق النجاح بالرسالة كاملة.

قد يكون المقطع الأخير وحده أقصر من ميزانية التعليق الذي يحمله، فيكمله البرنامج بأصفار. ويتيح طول الحمولة $L$ والطول الثابت لفحص السلامة حذف هذه الأصفار بعد جمع التتابع. وتشغل الأصفار جزءًا من الاختيارات المتاحة، دون أن تضيف محتوى للرسالة أو معلومات تحقق. وفي الإعداد الافتراضي على مستوى المنشور، تستخدم التعليقات القوائم المحضّرة المتفق عليها؛ فلا تصبح التعليقات الحاملة الجديدة اختيارات سرية إضافية في تلك القوائم الأصلية.

\subsubsection{جمع الرسالة لدى المُستقبِل وشروط رفضها}
يتلقى المُستقبِل المدخلات المشتركة وقائمة مرتبة بمراجع التعليقات الحاملة. ولا يكتشف هذا التنفيذ مجموعة الرسالة كاملة تلقائيًا من المنصة العامة. يحدد كل تعليق مشار إليه، ويستخرج موضع رده وموضوعه، ويحوّل الاختيارين إلى مقاطع بتات مرشحة. ثم يجمع المقاطع بالترتيب المزود ويقرأ التتابع الناتج. ويرفض الرسالة عند تعذر قراءة أحد التعليقات، أو عدم اتساق العدد أو الطول، أو إخفاق \EN{CRC}، أو وجود بتات غير صفرية في موضع الإكمال. وإذا لم تسمح التفسيرات المرشحة بتحديد حمولة واحدة قابلة للاستخراج، يبلغ عن الإخفاق. وبعد قبول الحمولة فقط، يفك الضغط ويعكس تحويل الحماية المتفق عليه.

يُقاس النجاح على مستوى الرسالة كاملة. فلا تكفي قراءة معظم التعليقات إذا منعت الاختيارات المتبقية إعادة بناء الرسالة. ولا يوفر حقل العدد تصحيحًا للأخطاء أو إجراءً لإعادة الإرسال. لذلك يجب أن ترافق كل نتيجة للتعليقات المتعددة شروط تجميع التعليقات وترتيبها ونشرها واستخراج الرسالة منها.

\subsubsection{حساب كلفة المثال ذي التعليقات السبعة}
المثال المسجّل في القسم~\ref{sec:multicomment-example} رسالة واحدة تحملها \NUM{\McCount} تعليقات. تضم الحمولة المحضّرة \NUM{\McCompressedBits} بتًا. ويحتاج العدد $N=\McCount$ إلى \NUM{\McCountBits} بتات بترميز العدد، ويحتاج الطول $L=\McCompressedBits$ إلى \NUM{\McLengthBits} بتات. لذلك يبلغ مجموع بتات التحكم:
\[
\McMarkerBits+\McVersionBits+\McCountBits+\McLengthBits=\McControlBits.
\]
ويضيف فحص \EN{CRC} المنفصل \NUM{\McIntegrityBits} بتًا، فتصبح كلفة النقل الإضافية \NUM{64} بتًا للرسالة كلها. ومع بتات الحمولة الـ\NUM{\McCompressedBits}، يصبح تتابع النقل بطول \NUM{\McStreamBits} بتًا. وتتيح التعليقات السبعة، بواقع \NUM{\McFrameBits} بتًا للتعليق، ميزانية قدرها \NUM{\McSlotBits} بتات، منها \NUM{\McPaddingBits} أصفار لإكمال التعليق الأخير. ولا تُضاف ترويسة من \NUM{32} بتًا إلى كل واحد من التعليقات السبعة.

في هذه الحمولة القصيرة بعد الضغط، تتجاوز معلومات النقل حجم تمثيل الرسالة. وتتوزع كلفة البادئة والإصدار والفحص الثابتة على الرسالة كاملة، بينما ينمو حقلا العدد والطول لوغاريتميًا. يفسر ذلك كلفة التنفيذ، دون افتراض أن ترويسته مثلى. ولا تُحمّل نتائج المقارنة بنمط التعليق الواحد كلفة هذا النقل المنفصل.
